\documentclass[12pt]{book}
\usepackage{precalc1}
\setcounter{chapter}{1}
\begin{document}

\chapter{Lines and Parabolas}

This week studies the two simplest non-constant function shapes: lines and parabolas.
They are familiar from algebra, but the goal here is to see the geometry that the
formulas are recording.

A line records a fixed rate of change. Every step to the right raises or lowers the
output by a constant amount, and geometrically that means the right triangles riding
along the line are all similar. A parabola records a different kind of structure: it
is the curve associated with a focus, an axis, and a symmetry that makes parallel rays
converge to a single point. These two shapes become the basic models for constant
change and changing change.

\section{Lines as Families of Similar Triangles}

Slope is often introduced as a formula, but the reason the formula works is
geometric. Take any two points on a non-vertical line and drop the right triangle
between them: run horizontally, then rise vertically to meet the line. Choose a
different pair of points and build the same kind of triangle.

The two triangles have the same angles, because the line meets every horizontal at
one fixed angle. Therefore the triangles are similar. Similar triangles have
proportional sides, so the ratio of rise to run is identical no matter which pair of
points generated it. That common ratio is the slope.
\begin{definition}{Slope}
For two points $(x_1,y_1)$ and $(x_2,y_2)$ with $x_1\neq x_2$, the \emph{slope} of
the line through them is
\[
m=\frac{y_2-y_1}{x_2-x_1}=\frac{\text{rise}}{\text{run}}.
\]
\end{definition}

\begin{visual}
A line is a family of similar right triangles stacked along it. Each triangle has a
horizontal leg (the run) and a vertical leg (the rise), and the slope $m$ is the
fixed ratio rise$/$run shared by all of them. Doubling the run doubles the rise; the
shape of the triangle never changes, only its size.
\end{visual}

\begin{center}
\begin{tikzpicture}[scale=0.85]
  \draw[->] (-0.5,0) -- (6.5,0) node[right] {$x$};
  \draw[->] (0,-0.5) -- (0,5) node[above] {$y$};
  \draw[pcblue,thick,domain=-0.3:6.2] plot (\x,{0.6*\x+0.8});
  \node[pcblue] at (5.6,4.7) {\small $f(x)=mx+b$};
  % first triangle
  \draw[pcred,thick] (1,1.4) -- (2,1.4) -- (2,2.0) -- cycle;
  \node[pcred,below] at (1.5,1.4) {\scriptsize run};
  \node[pcred,right] at (2,1.7) {\scriptsize rise};
  % second, larger triangle
  \draw[pcgreen,thick] (3.5,2.9) -- (5.5,2.9) -- (5.5,4.1) -- cycle;
  \node[pcgreen,below] at (4.5,2.9) {\scriptsize $2\times$ run};
  \node[pcgreen,right] at (5.5,3.5) {\scriptsize $2\times$ rise};
  % intercept
  \fill[pcblue] (0,0.8) circle (2pt) node[left] {\small $b$};
\end{tikzpicture}
\end{center}

Once the slope is known, the line's tilt is known, but the line can still slide up
or down. The number $b$ records where that tilted line crosses the $y$-axis, the
output at input $0$. The slope sets the tilt; the intercept sets the height. Together
they fix the line completely, which is the content of the slope-intercept form.
\begin{keyidea}
The slope-intercept form of a line is
\[
f(x)=mx+b,
\]
where $m$ is the common rise-over-run ratio of every similar triangle along the line
and $b=f(0)$ is the height at which the line crosses the $y$-axis.
\end{keyidea}

\begin{example}[reading slope and intercept]
The line $f(x)=\tfrac34 x-2$ has slope $\tfrac34$: every run of $4$ produces a rise
of $3$. It crosses the $y$-axis at $(0,-2)$. From the intercept, stepping right $4$
and up $3$ reaches $(4,1)$, a second point on the line.
\end{example}

The coefficient $\tfrac34$ is therefore not just a number attached to $x$. It is the
shape of every rise-run triangle on the line.

\begin{pitfall}
Slope is rise over run, not run over rise. The line $f(x)=\tfrac34x-2$ climbs $3$ for
every $4$ across, a gentle rise. Inverting the ratio to $\tfrac43$ describes a
steeper line and a different function.
\end{pitfall}

\section{Two Points Determine a Line}

A point cloud of exactly two points (with different $x$-coordinates) determines one
line and only one. This is the linear version of fitting a model to data: there are
two parameters, $m$ and $b$, and two points give enough information to determine them.

The construction runs in three stages. First compute the slope from the two points.
Then use one of the points to anchor the line in place. Finally rearrange the result
to slope-intercept form, where the slope and intercept can be read directly.
The slope comes straight from the definition. With slope and one known point
$(x_1,y_1)$ in hand, every other point $(x,y)$ on the line must produce that same
slope against $(x_1,y_1)$:
\[
\frac{y-y_1}{x-x_1}=m.
\]
Clearing the denominator gives the point-slope form. This form is useful because it
keeps visible the two pieces of information we actually know: one point and one
slope.
\begin{keyidea}
The point-slope form of the line through $(x_1,y_1)$ with slope $m$ is
\[
y-y_1=m(x-x_1).
\]
Solving for $y$ converts it to slope-intercept form $y=mx+b$.
\end{keyidea}

\begin{example}[from two points to $f(x)=mx+b$]
Find the line through $(2,1)$ and $(6,4)$.

We will separate the work into slope, point-slope form, and slope-intercept form.
Keeping those stages separate makes the method easier to check.

\emph{Slope.} $\displaystyle m=\frac{4-1}{6-2}=\frac{3}{4}.$

\emph{Point-slope.} Using $(2,1)$,
\[
y-1=\tfrac34(x-2).
\]

\emph{Slope-intercept.} Expand and solve for $y$:
\[
y=\tfrac34x-\tfrac32+1=\tfrac34x-\tfrac12,
\qquad\text{so}\qquad f(x)=\tfrac34x-\tfrac12.
\]
Both original points check: $f(2)=\tfrac34\cdot2-\tfrac12=\tfrac32-\tfrac12=1$ and
$f(6)=\tfrac34\cdot6-\tfrac12=\tfrac92-\tfrac12=4$.
\end{example}

\begin{center}
\begin{tikzpicture}[scale=0.8]
  \draw[->] (-0.5,0) -- (7.5,0) node[right] {$x$};
  \draw[->] (0,-1) -- (0,5.5) node[above] {$y$};
  \foreach \x in {1,...,7} \draw (\x,0.08)--(\x,-0.08) node[below] {\scriptsize \x};
  \foreach \y in {1,...,5} \draw (0.08,\y)--(-0.08,\y) node[left] {\scriptsize \y};
  \draw[pcblue,thick,domain=-0.3:7.3] plot (\x,{0.75*\x-0.5});
  \fill[pcred] (2,1) circle (2.4pt) node[above left] {\small $(2,1)$};
  \fill[pcred] (6,4) circle (2.4pt) node[above left] {\small $(6,4)$};
  % similar triangle between the two points
  \draw[pcgreen,thick,dashed] (2,1) -- (6,1) -- (6,4);
  \node[pcgreen,below] at (4,1) {\scriptsize run $=4$};
  \node[pcgreen,right] at (6,2.5) {\scriptsize rise $=3$};
\end{tikzpicture}
\end{center}

\begin{pitfall}
The slope must use the coordinates in consistent order: the $y$-difference and the
$x$-difference both run from the same starting point. Subtracting $y_2-y_1$ over
$x_1-x_2$ flips the sign and tilts the line the wrong way.
\end{pitfall}

\begin{example}[a horizontal and a vertical case]
Through $(1,3)$ and $(5,3)$ the slope is $\tfrac{3-3}{5-1}=0$, giving $f(x)=3$, a
horizontal line. A zero slope means the output does not change as the input changes.

Through $(2,1)$ and $(2,7)$ the run is $2-2=0$, so the slope formula divides by
zero. The relation $x=2$ is a vertical line, which fails the vertical line test and
so is not a function. This is not a line with an unusual slope; it is a relation that
cannot be written as $y=f(x)$.
\end{example}

\section{The Parabola as a Focusing Curve}

A parabola is not merely a bent line. It is the next basic shape after a line, and
its geometry is richer. Algebraically, it comes from a quadratic function. Geometrically,
it is the curve defined by a focusing property: every ray traveling parallel to its
axis reflects off the curve and passes through a single point, the focus.

This is why headlight reflectors, satellite dishes, and solar collectors are built
using parabolic geometry. The same curve can concentrate parallel incoming energy at
one point, or, run in reverse, send energy from the focus out as a parallel beam.
\begin{visual}
Parallel rays strike the inside of a parabola and all bounce to the focus. No other
curve collects parallel rays to a single point. The algebraic parabola $y=ax^2$ and
the optical reflector are the same object viewed two ways.
\end{visual}

\begin{center}
\begin{tikzpicture}[scale=1.0]
  % parabola y = 0.5 x^2, focus at (0, 0.5)
  \draw[pcblue,thick,domain=-2.2:2.2,smooth,samples=80] plot (\x,{0.5*\x*\x});
  \fill[pcred] (0,0.5) circle (2pt) node[right] {\small focus};
  % incoming parallel rays
  \foreach \x in {-1.8,-1.2,-0.6,0.6,1.2,1.8}{
    \draw[pcgreen,-{Stealth[length=2mm]}] (\x,3.2) -- (\x,{0.5*\x*\x+0.03});
    \draw[pcgreen] ({\x},{0.5*\x*\x}) -- (0,0.5);
  }
  \node[pcgreen] at (2.3,3.0) {\scriptsize parallel rays};
  \draw[dashed] (0,0) -- (0,3.4) node[above] {\scriptsize axis};
\end{tikzpicture}
\end{center}

The focusing property has an equivalent point-and-line description. This definition
looks different from the optical picture, but it describes the same curve and gives a
clean way to construct it from distances.
\begin{definition}{Parabola}
Fix a point $F$ (the \emph{focus}) and a line $\ell$ not through $F$ (the
\emph{directrix}). The \emph{parabola} is the set of all points equidistant from $F$
and $\ell$.
\end{definition}

So a point and a line determine a parabola through the equidistance condition, just
as two points determine a line. A point $P$ lies on the parabola exactly when its
distance to the focus equals its perpendicular distance to the directrix. The vertex
sits halfway between focus and directrix, and the axis is the perpendicular line from
the focus to the directrix.
\begin{center}
\begin{tikzpicture}[scale=1.1]
  \draw[pcblue,thick,domain=-2:2,smooth,samples=80] plot (\x,{0.5*\x*\x});
  \fill[pcred] (0,0.5) circle (2pt) node[right] {\small $F$};
  \draw[pcpurple,thick] (-2.3,-0.5) -- (2.3,-0.5) node[right] {\small $\ell$};
  % a point on the parabola with equal distances
  \coordinate (P) at (1.4,0.98);
  \fill[pcgreen] (P) circle (2pt) node[above right] {\small $P$};
  \draw[pcgreen,dashed] (P) -- (0,0.5) node[midway,above] {\scriptsize $d_1$};
  \draw[pcgreen,dashed] (P) -- (1.4,-0.5) node[midway,right] {\scriptsize $d_2$};
  \node at (3.0,0.6) {\scriptsize $d_1=d_2$};
  \fill (0,0) circle (1.4pt) node[below left] {\scriptsize vertex};
\end{tikzpicture}
\end{center}

\section{Three Points Determine a Parabola}

A vertical-axis parabola is the graph of a quadratic function
\[
f(x)=ax^2+bx+c,\qquad a\neq 0.
\]
The three numbers $a$, $b$, and $c$ control the parabola, so three pieces of
independent information are needed to determine them.

Three points with distinct $x$-coordinates give three equations. Solving that system
fixes the three parameters uniquely. In this sense, three points determine a
vertical-axis parabola in the same way that two points determine a line.
\begin{keyidea}
Two points (distinct $x$) determine a unique line, $f(x)=mx+b$, with two parameters.
Three points (distinct $x$) determine a unique vertical-axis parabola,
$f(x)=ax^2+bx+c$, with three parameters. The count of points matches the count of
parameters.
\end{keyidea}

\begin{example}[fitting a parabola to three points]
Find the quadratic through $(0,1)$, $(1,0)$, and $(2,3)$.

The unknowns are $a$, $b$, and $c$. Each point gives one equation by substituting its
$x$- and $y$-coordinates into $f(x)=ax^2+bx+c$:
\[
\begin{aligned}
(0,1):&\quad c=1,\\
(1,0):&\quad a+b+c=0,\\
(2,3):&\quad 4a+2b+c=3.
\end{aligned}
\]
From the first, $c=1$. The second gives $a+b=-1$. The third gives $4a+2b=2$, i.e.
$2a+b=1$. Subtracting, $a=2$, then $b=-3$. The quadratic is
\[
f(x)=2x^2-3x+1.
\]
Check $(2,3)$: $2\cdot4-6+1=3$.
\end{example}

This example is also a reminder that fitting a quadratic is a linear algebra problem
in disguise. The equations are linear in the unknown parameters $a$, $b$, and $c$,
even though the resulting function is quadratic in $x$.

\begin{pitfall}
The three points must have distinct $x$-coordinates. Two points sharing an
$x$-value would force two outputs at one input, violating the function condition,
and the system would have no solution.
\end{pitfall}

\section{Three Forms of a Quadratic}

A quadratic function can be written three common ways. The forms describe the same
parabola, but each one makes different information easy to see. Part of learning
quadratics is learning which form to use for the question being asked.
\begin{center}
\renewcommand{\arraystretch}{1.6}
\begin{tabular}{@{}lll@{}}
\toprule
\textbf{Form} & \textbf{Expression} & \textbf{What it shows at a glance}\\
\midrule
Standard      & $f(x)=ax^2+bx+c$        & $y$-intercept $c$; leading behavior $a$\\
Factored      & $f(x)=a(x-r_1)(x-r_2)$  & the roots (x-intercepts) $r_1,r_2$\\
Vertex        & $f(x)=a(x-h)^2+k$       & the vertex $(h,k)$\\
\bottomrule
\end{tabular}
\end{center}

In all three forms the same $a$ controls the opening: upward for $a>0$, downward for
$a<0$, and steeper as $|a|$ grows. The vertex form is the transformation form from
Week 1 with parent $f(x)=x^2$: $h$ shifts horizontally, $k$ shifts vertically, and
$a$ scales and possibly reflects. The same parabola is being viewed through three
different algebraic windows.
\begin{center}
\begin{tikzpicture}[scale=0.7]
  \draw[->] (-3.3,0) -- (3.3,0) node[right] {$x$};
  \draw[->] (0,-1.3) -- (0,4.5) node[above] {$y$};
  \draw[pcblue,thick,domain=-2.1:2.1,smooth] plot (\x,{(\x)*(\x)-0});
  % vertex form parabola: a(x-h)^2+k, shifted
  \draw[pcred,thick,domain=-0.6:3.0,smooth] plot (\x,{0.8*(\x-1.2)*(\x-1.2)+0.3});
  \fill[pcred] (1.2,0.3) circle (2.2pt) node[below right] {\scriptsize vertex $(h,k)$};
  \node[pcblue] at (-2.4,3.7) {\scriptsize $x^2$};
  \node[pcred] at (3.0,3.5) {\scriptsize $a(x-h)^2+k$};
\end{tikzpicture}
\end{center}

\subsection{Standard to vertex: completing the square}

Completing the square rewrites $ax^2+bx+c$ as $a(x-h)^2+k$ by reconstructing the
perfect square hidden in the first two terms. The goal is not to change the function,
but to reveal its vertex.

The procedure is systematic: factor $a$ from the $x$-terms, add and subtract the
square that completes the bracket, then collect the constant terms.
\begin{example}[completing the square]
Convert $f(x)=2x^2-12x+19$ to vertex form.

Factor $2$ from the variable terms:
\[
f(x)=2\big(x^2-6x\big)+19.
\]
The square completing $x^2-6x$ uses half of $-6$, squared, which is $9$:
\[
f(x)=2\big(x^2-6x+9-9\big)+19=2\big((x-3)^2-9\big)+19.
\]
Distribute and collect:
\[
f(x)=2(x-3)^2-18+19=2(x-3)^2+1.
\]
The vertex is $(3,1)$ and the parabola opens upward, stretched by $2$.
\end{example}

The completed square has moved the important geometry to the surface. In
$2(x-3)^2+1$, the numbers $3$ and $1$ are no longer hidden inside the expanded
quadratic; they are the coordinates of the vertex.

The vertex coordinates can also be read off a formula that completing the square
produces in general:
\[
h=-\frac{b}{2a},\qquad k=f(h).
\]

\begin{example}[vertex by formula]
For $f(x)=2x^2-12x+19$, $h=-\tfrac{-12}{2\cdot2}=3$ and $k=f(3)=18-36+19=1$. The
vertex is $(3,1)$, agreeing with the completed-square result.
\end{example}

\subsection{Standard to factored: the roots}

The factored form exposes the $x$-intercepts, the inputs where $f(x)=0$. These are
also called roots or zeros. They come from factoring directly when possible, or from
the quadratic formula in general:
\[
r_{1,2}=\frac{-b\pm\sqrt{b^2-4ac}}{2a}.
\]
The quantity under the root, $b^2-4ac$, is the \emph{discriminant}. It records how
many real roots exist before we even compute them: two when positive, one (a doubled
root at the vertex) when zero, and none when negative. Geometrically, it tells us
whether the parabola crosses the $x$-axis twice, touches it once, or misses it.
\begin{example}[standard to factored]
Factor $f(x)=x^2-x-6$. Two numbers multiplying to $-6$ and adding to $-1$ are $-3$
and $2$, so
\[
f(x)=(x-3)(x+2),
\]
with roots $r_1=3$ and $r_2=-2$. The discriminant is $(-1)^2-4(1)(-6)=25>0$,
confirming two real roots.
\end{example}

\begin{example}[roots by the quadratic formula]
For $f(x)=2x^2-3x+1$, the discriminant is $9-8=1$, so
\[
r=\frac{3\pm1}{4}=1\ \text{or}\ \tfrac12.
\]
The factored form is $f(x)=2(x-1)\!\left(x-\tfrac12\right)$. Expanding recovers
$2x^2-3x+1$.
\end{example}

\subsection{Relations among the parameters}

The three forms are not separate facts to memorize. They share structure, and a few
relations connect their parameters directly. These relations are useful because they
let us move from one piece of information to another without always expanding or
factoring from scratch.
\begin{keyidea}
For $f(x)=ax^2+bx+c=a(x-r_1)(x-r_2)=a(x-h)^2+k$:
\begin{itemize}\setlength{\itemsep}{2pt}
  \item the vertex sits at $h=-\dfrac{b}{2a}$, midway between the roots:
        $h=\dfrac{r_1+r_2}{2}$;
  \item the product of the roots is $r_1r_2=\dfrac{c}{a}$ and their sum is
        $r_1+r_2=-\dfrac{b}{a}$;
  \item the $y$-intercept is $c=f(0)=a\,r_1r_2=a h^2+k$.
\end{itemize}
\end{keyidea}

\begin{example}[reading the vertex from the roots]
The parabola $f(x)=(x-3)(x+2)$ has roots $3$ and $-2$. Their midpoint is
$h=\tfrac{3+(-2)}{2}=\tfrac12$, so the axis of symmetry is $x=\tfrac12$ and the
vertex lies there. Then $k=f\!\left(\tfrac12\right)=\left(\tfrac12-3\right)
\left(\tfrac12+2\right)=\left(-\tfrac52\right)\!\left(\tfrac52\right)=-\tfrac{25}{4}$.
The vertex is $\left(\tfrac12,-\tfrac{25}{4}\right)$.
\end{example}

The averaging step works because the two roots lie symmetrically on either side of
the axis. Once the axis is known, evaluating the function there gives the height of
the vertex.

\begin{visual}
A parabola is symmetric about the vertical line through its vertex. The two roots sit
at equal distances on either side of that axis, so their midpoint is the vertex's
$x$-coordinate. Reading the vertex from the roots is just averaging.
\end{visual}

\begin{center}
\begin{tikzpicture}[scale=0.8]
  \draw[->] (-3.3,0) -- (4.3,0) node[right] {$x$};
  \draw[->] (0,-2.2) -- (0,3.2) node[above] {$y$};
  \draw[pcblue,thick,domain=-2.7:3.7,smooth,samples=80]
    plot (\x,{0.5*(\x-3)*(\x+2)});
  \fill[pcred] (3,0) circle (2pt) node[above right] {\scriptsize $r_1=3$};
  \fill[pcred] (-2,0) circle (2pt) node[above left] {\scriptsize $r_2=-2$};
  \draw[pcpurple,dashed] (0.5,-2.2) -- (0.5,2.0);
  \node[pcpurple] at (1.5,2.0) {\scriptsize axis $x=\tfrac12$};
  \fill[pcgreen] (0.5,-1.5625) circle (2pt) node[below right] {\scriptsize vertex};
\end{tikzpicture}
\end{center}

\section{Exercises}

\begin{exercises}
\item Explain why the slope of a line is the same regardless of which two points are
used to compute it, referring to similar triangles in your answer.

\item Find the equation of the line through $(-1,2)$ and $(3,-6)$. Carry it through
slope, point-slope, and slope-intercept form, and verify with both points.

\item A line passes through $(4,1)$ with slope $-\tfrac23$. Write its point-slope and
slope-intercept forms, and find its $x$- and $y$-intercepts.

\item Decide whether each relation is a function and, if so, give its slope:
(a) the line through $(2,5)$ and $(2,-1)$; (b) the line through $(0,4)$ and $(7,4)$;
(c) the line through $(1,1)$ and $(3,5)$.

\item State the focus-directrix definition of a parabola and explain, in terms of
equidistance, why the vertex lies halfway between focus and directrix.

\item Describe the focusing property of a parabola and name two devices that exploit
it. Explain which direction the rays travel in each device.

\item Find the quadratic $f(x)=ax^2+bx+c$ through $(-1,6)$, $(0,1)$, and $(2,3)$ by
setting up and solving the linear system.

\item Convert $f(x)=3x^2+12x+7$ to vertex form by completing the square, and confirm
the vertex with the formula $h=-\tfrac{b}{2a}$.

\item For $f(x)=x^2-2x-15$, find the roots by factoring, the vertex by averaging the
roots, and the $y$-intercept. Sketch the parabola with all three labeled.

\item Use the discriminant to determine the number of real roots of each:
(a) $x^2-4x+4$; (b) $2x^2+x+3$; (c) $x^2-5x+2$. For any with real roots, give them.

\item A quadratic has roots $-4$ and $2$ and passes through $(0,-16)$. Determine $a$,
then write the function in all three forms.

\item Show, using $h=-\tfrac{b}{2a}$ and the root-sum relation $r_1+r_2=-\tfrac{b}{a}$,
that the vertex's $x$-coordinate equals the average of the two roots. Explain the
connection to the axis of symmetry.
\end{exercises}

\end{document}
